% Ad Atticum 9.11 (ad-atticum-09-11)
% Source: https://raw.githubusercontent.com/PerseusDL/canonical-latinLit/master/data/phi0474/phi057/phi0474.phi057.perseus-lat1.xml
% Fetched by scripts/fetch_latin.py

% Dateline (Perseus): Scr. in Formiano xiii K. Apr. a. 705(49).

\ciceroLetterOpener{CICERO ATTICO salutem}

\ciceroSection{1} Lentulum nostrum scis Puteolis esse? quod cum e viatore quodam esset auditum qui se diceret eum in Appia, cum is paulum lecticam aperuisset, cognosse, etsi vix veri simile est , misi tamen Puteolos pueros qui pervestigarent et ad eum litteras. inventus est vix in hortis suis se occultans litterasque mihi remisit mirifice gratias agens Caesari; de suo autem consilio C. Caesio mandata ad me dedisse. eum ego hodie exspectabam, id est xiii K. Aprilis.

\ciceroSection{2} venit etiam ad me Matius Quinquatribus, homo me hercule, ut mihi visus est, temperatus et prudens; existimatus quidem est semper auctor oti. quam ille hoc non probare mihi quidem visus est, quam illam νέκυιαν , ut tu appellas, timere! huic ego in multo sermone epistulam ad me Caesaris ostendi, eam cuius exemplum ad te antea misi, rogavique ut interpretaretur quid esset quod ille scriberet consilio meo se uti velle, gratia, dignitate, ope rerum omnium. respondit se non dubitare quin et opem et gratiam meam ille ad pacificationem quaereret. utinam aliquod in hac miseria rei publicae πολιτικὸν opus efficere et navare mihi liceat! Matius quidem et illum in ea sententia esse confidebat et se auctorem fore pollicebatur. pridie autem apud me Crassipes fuerat qui se pridie Non. Martias Brundisio profectum atque ibi Pompeium reliquisse dicebat, quod etiam qui viii Idus illinc profecti erant nuntiabant; illa vero omnes in quibus etiam Crassipes qui pro prudentia potuit attendere, sermones minacis, inimicos optimatium, municipiorum hostis, meras proscriptiones, meros Sullas; quae Lucceium loqui, quae totam Graeciam, quae vero Theophanem!

\ciceroSection{4} et tamen omnis spes salutis in illis est et ego excubo animo nec partem ullam capio quietis et, ut has pestis effugiam, cum dissimillimis nostri esse cupio! quid enim tu illic Scipionem, quid Faustum, quid Libonem praetermissurum sceleris putas quorum creditores convenire dicuntur? quid eos autem, cum vicerint, in civis effecturos? quam vero μικροψυχίαν Gnaei nostri esse? nuntiant Aegyptum et Arabiam εὐδαίμονα et Μεσοποταμίαν cogitare, iam Hispaniam abiecisse. monstra narrant; quae falsa esse possunt sed certe et haec perdita sunt et illa non salutaria. tuas litteras iam desidero. post fugam nostram numquam iam tantum earum intervallum fuit. misi ad te exemplum litterarum mearum ad Caesarem quibus me aliquid profecturum puto.
